\section{问题提出的背景}

\subsection{全球变暖背景下的冰冻圈观测需求}
环境与气候的快速变化已成为 21 世纪人类社会面临的重要挑战之一，冰冻圈（Cryosphere）对气候变化高度敏感，其变化会通过反照率反馈、淡水输入与海平面变化等途径影响气候系统与人类社会可持续发展\cite{IPCC2019SROCC,Lenton2019}. IPCC《海洋与冰冻圈特别报告》指出，海洋与冰冻圈正在发生显著且持续的变化，亟需开展长期、稳定、可追溯的观测与协同行动\cite{IPCC2019SROCC}。在此背景下，极区（尤其北极）作为全球变暖响应最敏感的区域之一，其增暖与海冰变化可能通过大气环流和天气气候异常等机制影响中纬度地区\cite{Cohen2014}。因此，获取长期、连续、空间覆盖广且精度可靠的极区冰雪关键参数数据，是认识冰冻圈变化规律与评估气候风险的重要基础。

\subsection{海冰积雪及雪深的重要性}
北极海冰表面积雪通过高反照率与热绝缘效应显著调制海冰的生长与消融过程：一方面，积雪可有效反射太阳短波辐射；另一方面，雪层较低的热导率会改变海冰与大气的热通量交换，从而影响海冰热力学过程与季节演变。雪深（snow depth）是表征海冰积雪状态的关键参数之一，也是海冰厚度反演中的重要输入变量；雪深的不确定性会直接传播到海冰厚度与体积估计中，从而影响极区航道、海冰—海洋模型与气候评估等应用。

然而，北极自然环境恶劣、空间尺度巨大，地面与野外观测难以满足“大范围、长时序、高分辨率”的数据需求。遥感技术因其非接触、广覆盖与可重复观测的优势，成为获取海冰与积雪参数的核心手段。但现有雪深遥感方案仍面临分辨率与精度之间的矛盾：例如被动微波产品覆盖广但分辨率较粗，且在深雪、复杂地形和变湿雪条件下不确定性较大；机载雷达/激光雷达精度高但成本高、覆盖有限。

\subsection{ICESat-2 多次散射信息为雪深反演提供新方法}
ICESat-2 搭载的 ATLAS 为 532\,nm 光子计数激光高度计，在雪覆盖区域的沿轨观测中，除雪面主峰回波外，常可在垂向回波剖面中观测到由雪体内部多次散射形成的“次表层长尾”信号。该长尾本质上编码了光子在雪体内传播的回散射路径长度分布信息，为不依赖雪前/雪后差分、也不需要参考 DEM 的雪深反演提供了新的信息源\cite{Hu2022,Lu2022Unc}。

Hu 等（2022）将辐射传输理论与蒙特卡洛模拟结合，指出多次散射回波对应的路径长度分布可用 Gamma 分布近似刻画，并给出了雪深与路径长度分布低阶矩之间的简洁解析关系：例如平均路径长度满足
\begin{equation}
  \langle L \rangle \approx 2H,
\end{equation}
二阶、三阶矩与雪深及漫散射系数 \(k_{sd}\) 之间也存在幂律关系，从而形成“由矩到雪深”的物理统计反演框架\cite{Hu2022}。
Lu 等（2022）进一步面向工程应用系统讨论了误差来源与质量控制：仪器冲激响应与 afterpulse 伪影、大气前向散射、雪表面粗糙度/坡度导致的脉冲展宽，以及日间太阳背景噪声与探测器暗噪声等，均会改变剖面形态并引入雪深偏差\cite{Lu2022Unc}。这些研究共同表明：ICESat-2 的多次散射信号具备可靠的雪深反演潜力，但要形成可复现、可推广的产品化流程，仍需系统化的算法实现、参数稳健性检验与不确定性量化。

\subsubsection{项目提出的原因}
雪深探测的具体方法主要包括光学卫星遥感观测、被动微波遥感观测、主动微波观测以及激光雷达高度计观测。光学遥感观测虽然能够提供高分辨率的积雪覆盖信息 \cite{Shao2025}，但记录的辐射量亮度值与雪深关系微弱，难以精确估算雪深，且易受云雾等大气因素影响。

与可见光和近红外相比，被动微波不依赖于太阳光的存在，且由于气溶胶粒子的波长一般比微波的波长小（微波波长范围一般在 1mm 至 1m 之间），微波能够轻易穿透云、雾、沙尘等组成的气溶胶粒子。因此，大气因素对被动微波观测方法的影响较小，使其在不同天气条件下能够提供相对稳定的雪深测量结果。利用不同频率的微波辐射与积雪相互作用可以建立与雪深、雪水当量等相关的模型，反演全天候、全天时的雪深信息，是现有大尺度雪深估算最常用的方法 \cite{Cao2001}。但由于被动微波亮温数据空间分辨率较粗，在混合像元和复杂地形下雪深反演存在较大的不确定性，难以反演区域尺度的雪深变化，因此被动微波遥感技术仅适用于大范围的雪深估算与分析。此外，被动微波辐射信号对雪湿度、密度和颗粒大小以及海冰浓度、类型和变形程度等参数非常敏感，容易导致反演的雪深存在偏差 \cite{Singh2012}。

相较于被动微波，主动微波雪深观测因其高空间分辨率和对积雪特性参数的敏感性，自 1980 年以来被广泛用于雪深监测。主动微波遥感通常利用合成孔径雷达（SAR）、干涉合成孔径雷达（InSAR）以及差分合成孔径雷达干涉测量（D-InSAR）技术实现。SAR 技术具有亚米级的空间分辨率，适用于对大范围区域的雪深监测 \cite{Bulovic2025}，能够提供详细的雪深信息，并且对于地形和植被等的遮挡具有较强的穿透能力，从而有效测量较厚的积雪层。然而，SAR 传感器接收到的各后向散射项往往难以区分，导致雪深估算结果不够准确 \cite{Das2008}。InSAR 技术通过对 SAR 影像的相位信息进行干涉处理测量雪深，但相位变化不仅受雪深影响，还受积雪结构和表面粗糙度的影响，且去相干性受时空基线影响较大 \cite{Evans2014}。D-InSAR 技术通过提取降雪前后同一区域的差分干涉相位信息建立几何函数关系估算雪深，在区域尺度山区流域中具有较高精度，但其精度受干涉相干性、复杂地形、相位解缠及积雪介电常数等多种因素制约。

激光雷达（LiDAR）由于不受积雪特性和含水量的干扰，且具有较高的高程测量精度，在季节性雪深、海冰厚度及积雪面积测量方面具有广泛应用前景。相较于微波雷达，LiDAR 系统具有更高的分辨率和测距精度，适用于多种地表和气象条件 \cite{Kwok2017}。然而，在地形复杂的山区，特别是当坡度大于 30° 时，LiDAR 的垂直精度大大降低，且由于卫星升降轨道交叉点有限，难以获取足够的雪深数据 \cite{Moudry2022}。

目前，主流雪深探测方式多采用“多轨测量法”，即通过将有雪期和无雪期的地面高度相减实现。这种方法存在显著局限性：在常年冻土或海冰区域，由于缺少无雪基准数据，无法实现反演；对于敏感区域或打击目标区域，难以预先获取地表数据；且观测间隔通常达数月，无法保证期间地表未发生形变。因此，亟需开发一种通过单次轨迹探测即可实现雪深直接测量的手段。

\subsection{本研究的意义和目的}
在复杂地形、常年积雪与海冰等区域，传统依赖“有雪/无雪”地表高程差的雪深反演方法往往受限于无雪基准缺失、观测间隔长以及地表形变等因素，难以实现稳定、快速且可推广的雪深获取。ICESat-2 搭载的光子计数激光高度计在雪覆盖区域的垂向回波剖面中，除雪面主峰外，常可观测到由雪体内部多次散射形成的次表层长尾信号。该信号本质上记录了光子在雪体内传播的路径长度分布信息，为在单次轨迹观测条件下、无需雪前/雪后差分与参考 DEM 的雪深直接反演提供了新的物理观测量基础。这一观测特性为突破现有激光测高雪深反演对多时相数据的依赖，建立更具普适性的雪深获取方法提供了重要契机。

基于上述背景，本研究旨在系统挖掘 ICESat-2 多次散射信息在雪深反演中的物理意义与应用潜力，构建一套可复现、可推广且具有明确物理约束的雪深反演方法框架。研究将从光子数据出发，建立稳定的回波剖面构建与次表层信号提取流程，通过辐射传输理论与统计矩方法，将雪体多次散射路径长度分布与雪深之间建立可解释的定量联系，并在此基础上实现雪深反演。同时，将针对探测器冲激响应与 afterpulse 伪影、大气前向散射、地形坡度与表面粗糙度引起的脉冲展宽以及背景噪声等关键误差来源开展系统分析，形成相应的质量控制与不确定性表征方案。通过与机载、地面观测或权威雪深产品的对比验证，评估该方法在不同区域与雪况条件下的适用性与精度表现。最终，本研究期望为获取沿轨高空间分辨率、具备厘米量级精度潜力的雪深信息提供一种新的技术路径，为海冰厚度反演、陆地水文与雪过程模拟以及冰冻圈变化研究提供更可靠的观测支撑。